% Kapitel 20, satt ur lectures/L20/README.md och info/examination.md.

\chapter{Praktisk tentamen 2: CAN-moduler}
\label{c20:ch}

{\large\itshape Tre timmar, individuellt, vid datorn. Omfattar projektets moduler: bittajming, CRC,
bitstoppning, ramsekvensering, registersemantik och SPI-transaktionen.\par}

\begin{chapterbox}{I det här kapitlet}
\begin{itemize}
  \item Varför en individuell tentamen följer på ett grupparbete.
  \item Tentamens form: skrivuppgifter med utdelad testbänk, och läsuppgifter på papper.
  \item De fyra slagen av uppgifter: varianter, reparationer, ett enskilt register, läsning.
  \item Vad som är tillåtet att använda \textemdash{} inklusive gruppens egen projektkod.
  \item Hur man förbereder sig, vad tentamen mäter, och vart konstruktionen tar vägen efteråt.
\end{itemize}
\end{chapterbox}

Det här är kursens andra och sista examinationstillfälle. Som \chapref{c9:ch} innehåller det ingen ny
teori; kapitlet säger vad tentamen mäter, vilken form uppgifterna har och hur man förbereder sig.
Poäng och betygsgränser står i bilaga~\ref{app:project}.

Tentamen omfattar \textbf{projektets moduler}, alltså \chapref{c10:ch} till \ref{c19:ch}:
bittajmingen, CRC:n, bitstoppningen åt båda hållen, ramsekvenseringen, registersemantiken och
SPI-transaktionen. Del~\ref{part:vhdl}:s material är inte borta \textemdash{} varje modul i
projektet är byggd av den synkrona processmallen, en räknare, ett skiftregister och en Mooremaskin
\textemdash{} men tyngdpunkten ligger på del~\ref{part:can}.

% ------------------------------------------------------------------------------------------
\section{Varför en individuell tentamen på ett grupparbete}
\label{c20:sec:why}

Projektet bedöms i grupp och står för halva kursen. Den här tentamen mäter något annat: att var och
en förstår det gruppen byggde, inte bara den del man själv skrev.

Det är inte en misstro mot arbetsformen utan en följd av den. En grupp om fyra eller fem delar upp
elva moduler mellan sig, och den uppdelningen är precis vad som gör projektet genomförbart på tio
pass. Baksidan är lika förutsägbar: den som skrev \code{crc15} kan komma undan utan att någonsin ha
behövt förstå varför \code{bit_timer} har två pulser, och den som skrev registerbanken utan att ha
tänkt på vad stoppningen gör åt CRC:n. Uppgifterna hämtas därför ur \emph{alla} delar av projektet,
oavsett vem som skrev vilken modul.

\begin{aside}
Det finns en praktisk konsekvens av det här som är värd att ta på allvar långt före tentamen:
\textbf{kodgranskningen under projektet är förberedelsen}. En grupp som granskar varandras PR:er på
riktigt, i stället för att godkänna dem osedda, har redan läst allt som tentamen frågar om.
\Chapref{c15:ch}:s kodgranskningsseminarium är tillfället att öva det; det är också det tillfälle som
betalar sig här.
\end{aside}

% ------------------------------------------------------------------------------------------
\section{Upplägg}
\label{c20:sec:format}

Samma form som \chapref{c9:ch}: fristående uppgifter, oberoende av varandra och ordnade ungefär efter
stigande svårighet.

De flesta är \textbf{skrivuppgifter} \textemdash{} en utdelad, självkontrollerande testbänk och en
beskrivning av modulen den driver, med portordning, typer och beteende. Uppgiften är att skriva
modulen så att testbänken passerar.

Därutöver finns \textbf{läsuppgifter}: ett tidsdiagram eller en SPI-transaktion att tolka på papper.
De har ingen testbänk, av precis det skäl som gör dem värda att ställa \textemdash{} att läsa en
vågform är en färdighet i sig, och den examineras inte av att skriva VHDL. De rättas mot ett facit,
och delvis rätt svar ger delpoäng på samma sätt som delvis korrekt logik gör.

\subsection*{De fyra slagen av uppgifter}
\begin{description}[style=unboxed,leftmargin=0pt,itemsep=0.8ex]
  \item[Varianter.] En \code{bit_timer} med en annan sampelpunkt eller en annan bithastighet. En
    stoppningsdetektor som räknar ett annat antal. En CRC med ett annat polynom. Modulen är
    bekant; den enda punkt där den skiljer sig från projektets är den punkt som kräver att du
    förstått den.
  \item[Reparationer.] En given tillståndsmaskin eller ett givet skiftregister som är
    \emph{nästan} rätt. Hitta felet med hjälp av testbänkens utskrift och åtgärda det. Det är
    \chapref{c14:ch}:s och \ref{c15:ch}:s arbetssätt satt på prov: läs vad testbänken rapporterar,
    och gå bakåt därifrån till raden som orsakade det.
  \item[Ett enskilt register.] Ett register ur registerbanken, med sin sättnings- och
    nollställningsregel, isolerat från resten. \Chapref{c18:ch}:s skillnad mellan en puls och en
    klibbig nivå är hela uppgiften.
  \item[Läsning.] Ett tidsdiagram eller en SPI-transaktion att tolka: vad ligger på bussen, och vad
    borde ha legat där?
\end{description}

Delvis korrekt logik ger delpoäng, så lämna aldrig in en tom fil för att modulen inte blev klar;
detsamma gäller ett halvt uttolkat tidsdiagram. Rättningen läser dessutom koden, inte bara
testbänkens utfall: en modul som passerar genom att göra något annat än det som efterfrågades ger
inte full poäng.

% ------------------------------------------------------------------------------------------
\section{Hjälpmedel}
\label{c20:sec:aids}

\textbf{Tillåtet:}
\begin{itemize}
  \item Allt kursmaterial: den här boken, föreläsningarnas README och samtliga appendix.
    Registerkartan och protokollspecifikationen (bilaga~\ref{app:protocol}) hör förstås dit; flera
    uppgifter är skrivna mot dem.
  \item \textbf{Gruppens egen projektkod.} Ni byggde den; det vore konstigt att låtsas annat.
    Uppgifterna är skrivna så att avskrift inte hjälper \textemdash{} varianterna skiljer sig på
    just den punkt som kräver att man förstått modulen, och reparationerna har ingen motsvarighet i
    repot alls.
  \item Din egen dator med GHDL, och CircuitVerse om du vill rita en krets eller en vågform innan du
    skriver den. Det är ofta snabbare, inte långsammare.
\end{itemize}

\textbf{Inte tillåtet:} kommunikation med andra i någon form; språkmodeller och andra verktyg som
genererar eller förklarar VHDL åt dig; och sökning på nätet efter färdiga lösningar. Gränsen går
inte vid vem som trycker på tangenterna utan vid vems förståelse som prövas.

Den fullständiga och auktoritativa hjälpmedelslistan står i bilaga~\ref{app:project}; står något
annat här gäller den. Är du osäker på om något är tillåtet: fråga före tentan, inte under.

% ------------------------------------------------------------------------------------------
\section{Före tentamen}
\label{c20:sec:prep}

\begin{itemize}
  \item Gå igenom de moduler i projektet som \emph{någon annan} i gruppen skrev. Det är där luckorna
    finns, och det är det enda rådet i listan som är svårt att följa dagen innan.
  \item Kontrollfrågorna sist i \chapref{c12:ch} till \ref{c19:ch} är avsiktligt skrivna som
    repetitionsmaterial. Kan du svara på alla utan att slå upp något är du klar.
  \item Läs igenom de utdelade testbänkarna en gång till. De är det tydligaste som finns om exakt
    vad varje modul lovar, och en skrivuppgift är alltid lättare för den som vet hur en testbänk
    formulerar sina krav.
  \item Rita om \code{bit_timer}s två pulser och \code{can_controller}s tillståndsdiagram på ett
    papper, ur minnet. De två bilderna bär mer av projektet än någon enskild fil gör.
  \item Se till att GHDL fungerar på din maskin \emph{innan} du kommer. En trasig verktygskedja är
    inget skäl till förlängd skrivtid.
\end{itemize}

% ------------------------------------------------------------------------------------------
\section{Det här mäter tentamen}
\label{c20:sec:measures}

\begin{itemize}
  \item Att du kan omsätta en specifikation i ord till en korrekt modul, med rätt portordning.
  \item Att du förstår bittajmingen: varför sampelpunkten ligger där den ligger, och vilken puls som
    gör vad.
  \item Att du kan resonera om bitstoppning åt båda hållen \textemdash{} att sätta in och att ta
    bort är samma regel sedd från två håll, och båda sidorna måste hålla CRC:n utanför.
  \item Att du kan skilja en puls från en klibbig nivå, och veta vilket lager som ansvarar för
    vilket.
  \item Att du kan läsa en testbänks utskrift och hitta tillbaka till orsaken.
\end{itemize}

% ------------------------------------------------------------------------------------------
\section{Efter kursen}
\label{c20:sec:after}

Kontrollern går vidare: parallellklassen skriver C++-drivrutinen mot samma register över samma
SPI-transport, och noden följer med till CAN-labben i nästa kurs i inbyggda system. Det som håller
ihop de två halvorna är registerkartan och protokollspecifikationen i bilaga~\ref{app:protocol}
\textemdash{} inte koden på någondera sidan. Det är själva poängen med ett gränssnitt: två grupper
som aldrig läst varandras källkod kan ändå mötas på en buss, därför att specifikationen mellan dem är
skriven innan någon av dem började.

Över tjugo kapitel har boken gått från en grind på ett papper till en CAN-nod som en mikrokontroller
kan prata med. Varje mönster på vägen är allmängods i digital konstruktion och följer med oförändrat
in i vad du än bygger härnäst:

\begin{itemize}
  \item dubbelvippsynkroniseraren (\chapref{c4:ch}), som varje asynkron ingång passerar,
  \item den klockade processmallen (\chapref{c3:ch}), som varje sekventiell modul i boken är skriven
    med,
  \item Mooremaskinen (\chapref{c8:ch}), som en uppräknad typ och ett \code{case},
  \item skiftregistret och räknaren (\chapref{c6:ch}) som byggstenar snarare än som specialfall, och
  \item till sist uppdelningen i block med ett tydligt ansvar var (\chapref{c11:ch}), som är det enda
    skälet till att elva moduler går att skriva av fem personer parallellt och ändå passa ihop.
\end{itemize}

Vill du gå vidare från konstruktionen i den här boken \textemdash{} till en kontroller med felräknare
och ACK-kontroll, till verifiering bortom en handskriven testbänk, eller till nästa protokoll byggt
med samma metod \textemdash{} pekar bilaga~\ref{app:reading} ut vad som finns att läsa härnäst.
